% Appendix F --- further results on the smaller groups (Q-0312, review flags 4, 12 and 19;
% records/affine/move-q0312-section-4.py). Section 4.7 of Q-0309 ("Smaller groups") read as a side
% branch before the theorem on the general linear row; the statements it held move here verbatim,
% with their `% shared with blueprint` markers and their prose, each followed by its proof, moved
% from Appendix B. Section 4.9 is now a one-page summary of the smaller groups that points here and
% to Appendices G and H. prop:similarity-stable-order-families, whose clauses the general linear row
% spends, moved to section 4.6 instead.

\section{Further results on the smaller groups}\label{app:smaller-groups}

The statements of this appendix fill the rows of the smaller groups in Table~\ref{tab:landscape},
which \S\ref{sec:smaller-groups} summarizes. They show what a smaller group gains: families tied to
one frame, a ray column whose isotropic member no longer fixes the family, and a stable order whose
families lie strictly between their neighbours. The proofs close each subsection.

\subsection{Frames}

The affine image of a product family cascades along its own frame's axes, and only the Gaussian loses
its frame.

% shared with blueprint prop:frame-families
\begin{proposition}[frame families]\label{prop:frame-families}
Let $d \ge 1$ and, for $i = 1,\dots,d$, let $(\nu^{(i)}_{s,s'})_{0 \le s \le s'}$ be probability
measures on $\RR$ with $\nu^{(i)}_{s,s} = \delta_0$ and
$\nu^{(i)}_{s,s''} = \nu^{(i)}_{s',s''} * \nu^{(i)}_{s,s'}$ for $s \le s' \le s''$. For
$s \le s'$ in $[0,\infty)^d$, coordinatewise, put
$\mu^{\otimes}_{s,s'} := \nu^{(1)}_{s_1,s'_1} \otimes \dots \otimes \nu^{(d)}_{s_d,s'_d}$, and for
$A \in \GL(d)$ put $\mu_{A;s,s'} := A_*\mu^{\otimes}_{s,s'}$, the \emph{frame family}.
\begin{enumerate}
\item \emph{(Cascade in the frame, and covariance.)}
$\widehat\mu_{A;s,s'}(\omega) = \prod_{i}\widehat\nu^{(i)}_{s_i,s'_i}\bigl((A^{\tp}\omega)_i\bigr)$;
$\mu_{A;s,s} = \delta_0$ and $\mu_{A;s,s''} = \mu_{A;s',s''} * \mu_{A;s,s'}$ for
$s \le s' \le s''$; and $B_*\mu_{A;s,s'} = \mu_{BA;s,s'}$ for $B \in \GL(d)$.
\item \emph{(The Gaussian frame family.)} If $b \ge 0$ and $\nu^{(i)}_{s,s'}$ is the centred
Gaussian of variance $2b(s'^2 - s^2)$ for every $i$, then $\mu_{A;0,s}$ is the centred Gaussian of
covariance $2b\,A\diag(s_1^2,\dots,s_d^2)A^{\tp}$, and so depends on $(A,s)$ only through
$A\diag(s^2)A^{\tp}$.
\item \emph{(Only the Gaussian loses its frame.)} Let $d \ge 2$, and suppose $\mu_{A;0,s}$ depends
on $(A,s) \in \GL(d) \times (0,\infty)^d$ only through $A\diag(s^2)A^{\tp}$. Then there is a
$b \ge 0$ with $\nu^{(i)}_{0,t}$ the centred Gaussian of variance $2bt^2$ for every $i$ and every
$t > 0$ (the point mass $\delta_0$ when $b = 0$).
\end{enumerate}
\end{proposition}

Clauses (2) and (3) say that the frame drops out exactly when every factor is the same Gaussian,
with one constant $b$ for every axis. The core of (3) is Kelker's lemma (Remark~\ref{rem:kelker};
\sscite{kelker1970distribution}, Lemma~5, p.~423): with $A$ orthogonal and equal scales, the hypothesis
makes the product law of the factors rotation invariant.

The converse holds in a fixed frame, under weak hypotheses: covariance under the positive diagonal
matrices only, with no rotation and no reflection, and steps along one axis at a time.

% shared with blueprint prop:fixed-frame-families
\begin{proposition}[a fixed frame: the families are products]\label{prop:fixed-frame-families}
Let $d \ge 1$. For $s \in (0,\infty)^d$ write $D_s := \diag(s_1,\dots,s_d)$, and $\mathbf 1 :=
(1,\dots,1)$. Let $(\mu_s)_{s\in(0,\infty)^d}$ be probability measures on $\RR^d$ with
\begin{enumerate}
\item[(D1)] \emph{(covariance under the diagonal group)} $(D_\delta)_*\mu_s = \mu_{D_\delta s}$ for
all $\delta, s \in (0,\infty)^d$;
\item[(D2)] \emph{(a cascade along each axis)} for every $i \le d$ and every $\lambda > 1$ there is
a probability measure $\rho_{i,\lambda}$ on $\RR^d$ with
$\mu_{\mathbf 1 + (\lambda - 1)e_i} = \mu_{\mathbf 1} * \rho_{i,\lambda}$.
\end{enumerate}
Then there are self-decomposable laws $\nu_1, \dots, \nu_d$ on $\RR$
(in the sense of \ref{eq:self-decomposable} at $d = 1$) with
\[
  \mu_s = (s_1)_*\nu_1 \otimes \dots \otimes (s_d)_*\nu_d \qquad \text{for every } s \in (0,\infty)^d ,
\]
where $(s_i)_*$ is the image under $x \mapsto s_ix$ on $\RR$. Conversely, for any self-decomposable
laws $\nu_1, \dots, \nu_d$ on $\RR$ this family satisfies (D1) and (D2), and it cascades along
the coordinatewise order: $\mu_{s'} = \mu_s * (\rho^{(1)}_{s_1,s'_1} \otimes \dots \otimes \rho^{(d)}_{s_d,s'_d})$ for $s \le s'$,
with probability measures $\rho^{(i)}_{s_i,s'_i}$ on $\RR$.
\end{proposition}

So the non-Gaussian anisotropic families that cascade at a fixed orientation are exactly the
products of $d$ independent self-decomposable laws, one per axis, and the factors may differ. By
Proposition~\ref{prop:frame-families}(3) the non-Gaussian factors then fix the frame. A product family
gains anisotropy along its frame's axes as fast per step as the Gaussian does and cannot turn, while
the stable orbits turn, with a smaller gain per step (Proposition~\ref{prop:stable-step-cone}). Turning alone therefore does not
exclude non-Gaussian kernels. The proof of the product form uses the characteristic function of
Lemma~\ref{lem:charfn-one-on-subspace}.

\begin{proof}[Proof of Proposition~\ref{prop:frame-families}]
\emph{Clause (1).} The transform of a product measure is the product of the transforms, and
$\widehat{A_*\mu}(\omega) = \widehat\mu(A^{\tp}\omega)$. The cascade is coordinatewise multiplication
of transforms, and $(BA)_* = B_*A_*$ gives the covariance.

\emph{Clause (2).} By clause (1),
$\widehat\mu_{A;0,s}(\omega) = \exp[-b\sum_is_i^2(A^{\tp}\omega)_i^2] =
\exp[-b\langle\omega, A\diag(s^2)A^{\tp}\omega\rangle]$.

\emph{Clause (3).} For $Q \in \Orth(d)$, $Q\diag(\mathbf 1)Q^{\tp} = I$, so the hypothesis makes the
product law $\mu^\otimes_{0,\mathbf 1}$ rotation invariant: its transform is $f(|\omega|^2)$ with $f$
continuous, real and $f(0) = 1$, and each factor has the transform $t \mapsto f(t^2)$. As $d \ge 2$,
$\omega = \sqrt x\,e_1 + \sqrt y\,e_2$ gives $f(x+y) = f(x)f(y)$. Then $f > 0$, and $G := -\log f$ is
continuous, nonnegative and additive, so $G(u) = bu$ by the last paragraph of the proof of
Lemma~\ref{lem:single-ratio-cauchy}, the same Cauchy equation on which
Theorem~\ref{thm:affine-rigidity} rests. So each $\nu^{(i)}_{0,1}$ is the centred Gaussian of variance
$2b$, and $A = \diag(s)^{-1}$ in the hypothesis gives the variance $2bs_i^2$ at $s$.
\end{proof}

\begin{proof}[Proof of Proposition~\ref{prop:fixed-frame-families} (sketch)]
(D1) gives
$\mu_s = (D_s)_*\mu_{\mathbf 1}$; put $\varphi := \widehat\mu_{\mathbf 1}$.

\emph{A step along axis $i$ is a blur along axis $i$.} With $H := e_i^\perp$, (D2) gives
$\varphi(\lambda xe_i + \eta) = \varphi(xe_i + \eta)\,\widehat\rho_{i,\lambda}(xe_i + \eta)$. At $x = 0$,
$\widehat\rho_{i,\lambda} = 1$ on a neighbourhood of $0$ in $H$, so by
Lemma~\ref{lem:charfn-one-on-subspace}, in its neighbourhood form, $\widehat\rho_{i,\lambda}(xe_i + \eta)
= \widehat\rho_{i,\lambda}(xe_i)$. The same identity at $\eta = 0$ eliminates
$\widehat\rho_{i,\lambda}$:
\[
  \varphi(\lambda xe_i + \eta)\,\varphi(xe_i) = \varphi(xe_i + \eta)\,\varphi(\lambda xe_i)
  \qquad (x \in \RR,\ \eta \in H,\ \lambda > 1).
\]
For $x_0 \ne 0$ and $x$ strictly between $0$ and $x_0$ take $\lambda := x_0/x$. Letting $x \to 0$,
continuity and $\varphi(0) = 1$ give $\varphi(x_0e_i + \eta) = \varphi(\eta)\,\varphi(x_0e_i)$, which splits
off axis $i$.

\emph{The product.} Repeating over the axes, $\varphi(\omega) = \prod_i\varphi(\omega_ie_i)$, so
$\mu_{\mathbf 1}$ is the product of its marginals $\nu_i$. The image of (D2) under the $i$-th coordinate
is $(\lambda)_*\nu_i = \nu_i * \rho'$, which is self-decomposability on the line.

\emph{Converse.} $D_\delta$ acts coordinatewise, and self-decomposability of each factor gives the
increments, a convolution of product measures being the product of the convolutions. This is
Proposition~\ref{prop:frame-families}(1) in the standard frame.
\end{proof}

\subsection{The ray column}

Along rays a family of any group is a continuous covariant field of self-decomposable laws
(Proposition~\ref{prop:gw-ray-families}(2)), and the bookkeeping of Lemma~\ref{lem:gw-kernels-from-origin}
serves every cell below. For the similarities the parameter set is a $(d-1)$-parameter set of orbits,
and the isotropic member no longer determines the family.

% shared with blueprint prop:similarity-ray-families+21f2aac1312b
\begin{proposition}[similarity covariance along rays: the families, and the parameter set]\label{prop:similarity-ray-families}
Let $d \ge 1$, let $\mathrm S(d) := \{cQ : c > 0,\ Q \in \Orth(d)\}$ be the similarities, and let $W$ be
the ray order. For $\Sigma \in \PSD(d) \setminus \{0\}$ write $\sigma(\Sigma) \in \RR^d$ for its
eigenvalues in nonincreasing order divided by $\tr\Sigma$, and put
$\Delta_d := \{\sigma \in \RR^d : \sigma_1 \ge \dots \ge \sigma_d \ge 0,\ \sigma_1 + \dots + \sigma_d = 1\}$.
\begin{enumerate}
\item \emph{(The parameter set.)} $(\mathrm S(d), W)$ is a wedge datum. Two matrices
$\Sigma, \Sigma' \in \PSD(d) \setminus \{0\}$ lie in one orbit of $\mathrm S(d)$, acting by
$A\cdot\Sigma = A\Sigma A^{\tp}$, exactly when $\sigma(\Sigma) = \sigma(\Sigma')$, and each orbit contains
exactly one matrix $\diag(\sigma)$ with $\sigma \in \Delta_d$; so the orbits in $\PSD(d) \setminus \{0\}$
are indexed by $\Delta_d$, a simplex of dimension $d - 1$. The stabilizer of $\Sigma \ne 0$ in
$\mathrm S(d)$ is $\{Q \in \Orth(d) : Q\Sigma = \Sigma Q\}$, and it is $\Orth(d)$ exactly when $\Sigma$ is
a multiple of $I$.
\item \emph{(The families.)} The kernels from scale zero of the $(\mathrm S(d), W)$-cascade families of
kernels are the maps of Proposition~\ref{prop:gw-ray-families}(2) with $G = \mathrm S(d)$. In such a family every
$\mu_\Sigma$ is symmetric, $(-I)_*\mu_\Sigma = \mu_\Sigma$, and invariant under the stabilizer of
$\Sigma$, and $\mu_{cQ\diag(\sigma)Q^{\tp}} = (\sqrt c\,Q)_*\mu_{\diag(\sigma)}$ for $c > 0$, $Q \in \Orth(d)$
and $\sigma \in \Delta_d$: the family is determined by its kernels at the matrices $\diag(\sigma)$,
$\sigma \in \Delta_d$.
\item \emph{(The isotropic member does not determine the family.)} In such a family
$\mu_{sI} = (\sqrt s\,I)_*\mu_I$ for $s \ge 0$, and $\mu_I$ is rotation invariant, self-decomposable and
not $\delta_0$. Conversely, let $\mu_I$ be such a law. Then (i) $\Sigma \mapsto (\Sigma^{1/2})_*\mu_I$ and
(ii) $\Sigma \mapsto \nu_\Sigma := ((\det\Sigma)^{1/(2d)}I)_*\mu_I$ are the kernels from scale zero of
$(\mathrm S(d), W)$-cascade families with isotropic member $\mu_I$, and (ii) is also the family of kernels
from scale zero of an $(\mathrm S(d), \text{Loewner})$-cascade family. For $d \ge 2$ the two differ at
every $\Sigma \in \PD(d)$ that is not a multiple of $I$.
\item \emph{(The isotropic members of the similarity row.)} Let $W'$ be an $\mathrm S(d)$-invariant
partial order on $\PSD(d)$ that contains the rays and is contained in the Loewner order. The isotropic
members $\mu_I$ of the $(\mathrm S(d), W')$-cascade families of kernels are exactly the rotation-invariant
self-decomposable laws other than $\delta_0$.
\end{enumerate}
\end{proposition}

So the cell (similarities, rays) is large: a continuous field of symmetric self-decomposable laws over
the shapes, covariant under similarities. The kernel at a shape is constrained pointwise only by its
stabilizer, besides self-decomposability and continuity where eigenvalues meet, and it need not be
elliptical. The isotropic member is any rotation-invariant self-decomposable law other than $\delta_0$,
as in the general linear row, but for $d \ge 2$ it no longer fixes the family. Clause (4) says that the
isotropic members are the same in every column of the similarity row, the Loewner column included,
while in the general linear row the Loewner order leaves the Gaussian alone. Which choices of kernels at
the matrices $\diag(\sigma)$ extend to a continuous field is problem~\ref{prob:continuous-fields} of
\S\ref{sec:open-problems}.

For the positive diagonal group the parameter set is larger still: an orbit is one frame's own
shapes, and the stabilizer of $I$ is trivial, so $\mu_I$ need carry no symmetry.

% shared with blueprint prop:diagonal-ray-families+4ddf66022bc4
\begin{proposition}[diagonal covariance along rays: the families, and the parameter set]\label{prop:diagonal-ray-families}
Let $d \ge 1$, let $\mathrm D_+(d)$ be the group of diagonal matrices with positive diagonal entries, and
let $W$ be the ray order. For $\Sigma \in \PSD(d)$ write $\Delta(\Sigma) := \diag(\Sigma_{11}, \dots,
\Sigma_{dd})$, for $s \in (0,\infty)^d$ write $D_s := \diag(s_1, \dots, s_d)$, and put
$\mathrm{Corr}(d) := \{\Theta \in \PD(d) : \Theta_{11} = \dots = \Theta_{dd} = 1\}$.
\begin{enumerate}
\item \emph{(The parameter set.)} $(\mathrm D_+(d), W)$ is a wedge datum. Every $\Sigma \in \PD(d)$ has
the trivial stabilizer in $\mathrm D_+(d)$, and is $D\cdot \Theta = D\Theta D$ for exactly one pair
$D \in \mathrm D_+(d)$, $\Theta \in \mathrm{Corr}(d)$, namely $D = \Delta(\Sigma)^{1/2}$. So the orbits in
$\PD(d)$ are indexed by $\mathrm{Corr}(d)$, a nonempty open subset of the affine space of symmetric
matrices with unit diagonal, of dimension $d(d-1)/2$; the orbit of $I$ is the set of positive definite
diagonal matrices.
\item \emph{(The families on the frame's shapes.)} The kernels from scale zero of the
$(\mathrm D_+(d), W)$-cascade families of kernels are the maps of Proposition~\ref{prop:gw-ray-families}(2) with
$G = \mathrm D_+(d)$. On the orbit of $I$ they are $\mu_{D_s^2} = (D_s)_*\mu_I$, $s \in (0,\infty)^d$, with
$\mu_I$ self-decomposable and not $\delta_0$. Conversely, for every self-decomposable $\mu_I \ne \delta_0$
the map $\Sigma \mapsto (\Delta(\Sigma)^{1/2})_*\mu_I$ is the family of kernels from scale zero of a
$(\mathrm D_+(d), W)$-cascade family.
\item \emph{(The Loewner order: products.)} The kernels from scale zero of a
$(\mathrm D_+(d), \text{Loewner})$-cascade family of kernels are, on the orbit of $I$,
\[
  \mu_{D_s^2} = (s_1)_*\nu_1 \otimes \dots \otimes (s_d)_*\nu_d , \qquad s \in (0,\infty)^d ,
\]
with self-decomposable laws $\nu_1, \dots, \nu_d$ on $\RR$, not all $\delta_0$. Conversely, for such laws
the map $\Sigma \mapsto (\Delta(\Sigma)^{1/2})_*(\nu_1 \otimes \dots \otimes \nu_d)$ is the family of kernels
from scale zero of a $(\mathrm D_+(d), \text{Loewner})$-cascade family.
\item \emph{(Rays only: no product.)} Let $d \ge 2$ and $\Theta \in \mathrm{Corr}(d)$ not diagonal. The centred
Gaussian law $N(0,\Theta)$ is self-decomposable and is not of the form $\nu_1 \otimes \dots \otimes \nu_d$. So
the $(\mathrm D_+(d), W)$-family of (2) with $\mu_I = N(0,\Theta)$ is the restriction of no
$(\mathrm D_+(d), \text{Loewner})$-family.
\item \emph{(Off the frame's shapes.)} Let $d \ge 2$ and $\mu_I$ rotation invariant, self-decomposable and
not $\delta_0$. The maps $\Sigma \mapsto (\Delta(\Sigma)^{1/2})_*\mu_I$ and
$\Sigma \mapsto (\Sigma^{1/2})_*\mu_I$ are the kernels from scale zero of two
$(\mathrm D_+(d), W)$-cascade families with the same isotropic member, and they differ at every
$\Sigma \in \PD(d)$ that is not diagonal.
\end{enumerate}
\end{proposition}

So the cell (diagonal, rays) is not empty. On the orbit of $I$ it is Proposition~\ref{prop:ray-families}
without rotations: the frame images of one self-decomposable law, any such law. Clauses (3) and (4)
locate the product form of Proposition~\ref{prop:fixed-frame-families}. It is forced by the steps along
single axes, which the Loewner order has and the ray order lacks, and not by the group. Off the orbit of
$I$ the ray order links nothing across the correlation matrices. The continuous covariant fields
over $\mathrm{Corr}(d)$ are part of problem~\ref{prob:continuous-fields} of \S\ref{sec:open-problems}.

\begin{proof}[Proof of Proposition~\ref{prop:similarity-ray-families} (sketch)]
\emph{(1).} $\mathrm S(d)$ is closed in $\GL(d)$, as $AA^{\tp} = (AA^{\tp})_{11}\,I$ is a closed
condition and an invertible $A$ satisfying it is a similarity. A
similarity multiplies the eigenvalues and the trace by the same factor, so $\sigma$ is constant on
orbits, and the spectral theorem gives the converse. $\Delta_d$ is the convex hull of the affinely
independent points $k^{-1}(e_1 + \dots + e_k)$, $1 \le k \le d$. A similarity fixing $\Sigma \ne 0$ has $c = 1$ by
the traces, and a $\Sigma$ commuting with every orthogonal matrix is a multiple of $I$, by the
reflections and permutations of the proof of Lemma~\ref{lem:affine-gauge}.

\emph{(2)} is Proposition~\ref{prop:gw-ray-families}(2) with (b) of
Lemma~\ref{lem:gw-kernels-from-origin}(2) at $A = -I$, at the stabilizer, and at $A = \sqrt c\,Q$.

\emph{(3).} $\mu_I$ is invariant under the stabilizer $\Orth(d)$ of $I$, as in
Lemma~\ref{lem:parameter-space}(5). Family (i) is Proposition~\ref{prop:gw-ray-families}(3) restricted by
Lemma~\ref{lem:gw-kernels-from-origin}(1). For (ii), $\det$ is monotone along the Loewner order, so with
$\eta := (\det\Sigma)^{1/(2d)} \le \eta' := (\det\Sigma')^{1/(2d)}$ self-decomposability at
$\eta/\eta'$ gives the increment, and Lemma~\ref{lem:gw-kernels-from-origin}(3) gives the family.
For $d \ge 2$, if the two kernels agree at $\Sigma \in \PD(d)$, then with $t := (\det\Sigma)^{1/(2d)}$
the matrix $t^{-1}\Sigma^{1/2}$ fixes $\mu_I$, which is rotation invariant and not $\delta_0$, so it is
orthogonal by Lemma~\ref{lem:parameter-space}(2), and $\Sigma = t^2I$.
% CHANGED Q-0281 wave 2: blind review W2E-4: Lemma lem:parameter-space(2), as in the blueprint proof.

\emph{(4).} A family along $W'$ restricts to the rays, and the family (ii) of (3) restricts to $W'$,
which is contained in the Loewner order.
\end{proof}

\begin{proof}[Proof of Proposition~\ref{prop:diagonal-ray-families} (sketch)]
\emph{(1).} $(D\Sigma D)_{ii} = s_i^2\Sigma_{ii}$ forces $D = I$ on a stabilizer, and
$\Theta := \Delta(\Sigma)^{-1/2}\Sigma\Delta(\Sigma)^{-1/2}$ is the unique correlation matrix with
$\Delta(\Sigma)^{1/2}\cdot \Theta = \Sigma$. $\mathrm{Corr}(d)$ is the intersection of the affine space with the
open set of matrices whose quadratic form is positive on the compact unit sphere.

\emph{(2).} The first sentence is Proposition~\ref{prop:gw-ray-families}(2), and (b) of
Lemma~\ref{lem:gw-kernels-from-origin}(2) at $A = D_s$ gives $\mu_{D_s^2} = (D_s)_*\mu_I$. Conversely
$\Delta(D\Sigma D)^{1/2} = D\,\Delta(\Sigma)^{1/2}$ gives (b) for $\Sigma \mapsto (\Delta(\Sigma)^{1/2})_*\mu_I$,
and (a), (c) and self-decomposability are immediate.

\emph{(3).} On the orbit of $I$ the axis steps $I \le D^2_{\mathbf 1 + (\lambda-1)e_i}$ are Loewner steps,
so (D1) and (D2) of Proposition~\ref{prop:fixed-frame-families} hold and the product form follows.
Conversely $\Sigma \le \Sigma'$ gives $\Sigma_{ii} \le \Sigma'_{ii}$, and self-decomposability of each
factor gives a product increment.

\emph{(4).} $N(0,\Theta)$ is self-decomposable, and a product with unit diagonal would force $\Theta = I$.

\emph{(5).} If the two kernels agree at $\Sigma \in \PD(d)$, then with $D := \Delta(\Sigma)^{1/2}$ the
matrix $D^{-1}\Sigma^{1/2}$ fixes $\mu_I$, which is rotation invariant and not $\delta_0$, so it is
orthogonal by Lemma~\ref{lem:parameter-space}(2), and $\Sigma = D^2$ is diagonal.
% CHANGED Q-0281 wave 2: blind review W2E-4: Lemma lem:parameter-space(2), as in the blueprint proof.
\end{proof}

\subsection{The diagonal group along the stable order}

On the frame's shapes the stable order's steps are the diagonal exponents of the cone $\mathfrak k_\alpha$, so the
isotropic member must be operator-self-decomposable for each of them. The axis steps that force products
along the Loewner order are rank-deficient and absent.

% shared with blueprint prop:diagonal-stable-order-families
\begin{proposition}[diagonal covariance along the stable order: the families on the frame's shapes]\label{prop:diagonal-stable-order-families}
Let $d \ge 2$ and $0 < \alpha < 2$. Let $\mathrm D_+(d)$, $D_s$, $\Delta(\Sigma)$ and $\mathrm{Corr}(d)$ be as
in Proposition~\ref{prop:diagonal-ray-families}, let $\le_\alpha$ be the stable order on $\PSD(d)$ of
Proposition~\ref{prop:similarity-stable-order-families}(1), and use $B$-self-decomposability as in
Proposition~\ref{prop:stable-exponent-cone}. Write $\mathcal D_\alpha$ for the set of diagonal matrices $\diag(b)$,
$b \in \RR^d$, with $(d+\alpha)\min_ib_i \ge b_1 + \dots + b_d$; these are the diagonal matrices of the
cone $\mathfrak k_\alpha$ of Proposition~\ref{prop:stable-exponent-cone}.
\begin{enumerate}
\item \emph{(The datum.)} $(\mathrm D_+(d), \le_\alpha)$ is a wedge datum, and its orbits in $\PD(d)$ are
those of Proposition~\ref{prop:diagonal-ray-families}(1), indexed by $\mathrm{Corr}(d)$. On the orbit of $I$,
$D_s^2 \le_\alpha D_{s'}^2$ if and only if $u := (s'_1/s_1, \dots, s'_d/s_d)$ satisfies
$(\min_iu_i)^{d+\alpha} \ge u_1\cdots u_d$.
\item \emph{(On the frame's shapes.)} The kernels from scale zero of a
$(\mathrm D_+(d), \le_\alpha)$-cascade family of kernels are, on the orbit of $I$,
$\mu_{D_s^2} = (D_s)_*\mu_I$, $s \in (0,\infty)^d$, where $\mu_I \ne \delta_0$ is $B$-self-decomposable for
every $B \in \mathcal D_\alpha$; in particular $\mu_I$ is self-decomposable.
\item \emph{(Two sources of families.)} (a) For self-decomposable laws $\nu_1, \dots, \nu_d$ on $\RR$, not
all $\delta_0$, the map $\Sigma \mapsto (\Delta(\Sigma)^{1/2})_*(\nu_1 \otimes \dots \otimes \nu_d)$ is the
family of kernels from scale zero of a $(\mathrm D_+(d), \le_\alpha)$-cascade family of kernels.
(b) For a profile $k$ as in Proposition~\ref{prop:orbit-step-cone} with $R(\lambda) \ge \lambda^\alpha$ for every
$\lambda \ge 1$, and $\mu_I$ the law of that proposition's clause (1), so is $\Sigma \mapsto (\Sigma^{1/2})_*\mu_I$;
in particular for $\mu_I = S_\beta(I)$, $\alpha \le \beta < 2$.
\item \emph{(Strictly between its neighbours.)} (a) Let $|r| < 1$ with $r^2 > 4(1+\alpha)/(2+\alpha)^2$, and
let $\Theta_r \in \mathrm{Corr}(d)$ have the entries $(\Theta_r)_{12} = (\Theta_r)_{21} = r$ and agree with $I$ elsewhere.
The centred Gaussian law $N(0, \Theta_r)$ is self-decomposable and is not $B$-self-decomposable for
$B := I + \alpha\,e_1e_1^{\tp} \in \mathcal D_\alpha$. So the $(\mathrm D_+(d), \text{rays})$-family of
Proposition~\ref{prop:diagonal-ray-families}(2) with $\mu_I = N(0, \Theta_r)$ is the restriction of no
$(\mathrm D_+(d), \le_\alpha)$-family.
(b) $S_\alpha(I)$ is not of the form $\nu_1 \otimes \dots \otimes \nu_d$. So the family of (3)(b) with
$\mu_I = S_\alpha(I)$ is the restriction of no $(\mathrm D_+(d), \text{Loewner})$-family.
(c) Let $\sigma_\alpha$ be the law on $\RR$ with transform $e^{-|\tau|^\alpha}$ (Lemma~\ref{lem:elliptical-stable}(2)
at $d = 1$). The map $\Sigma \mapsto (\Sigma^{1/2})_*S_\alpha(I) * (\Delta(\Sigma)^{1/2})_*(\sigma_\alpha
\otimes \dots \otimes \sigma_\alpha)$ is the family of kernels from scale zero of a
$(\mathrm D_+(d), \le_\alpha)$-cascade family that is the restriction neither of a
$(\GL(d), \le_\alpha)$-family nor of a $(\mathrm D_+(d), \text{Loewner})$-family.
\end{enumerate}
\end{proposition}

So on the frame's shapes the cell (diagonal, stable order) lies strictly between its two neighbours in
the diagonal row: every restricted (diagonal, Loewner) family belongs to it, every family of it
restricts to the ray cell, neither inclusion is an equality, and the cell is not the union of the
restrictions of the (diagonal, Loewner) cell and the general linear cell. The condition on the isotropic
member sits between self-decomposability, which the ray cell requires, and the product form, which the
Loewner order's axis steps force.
% Round 3 (content review), E item 10: "its converse" said as what it asks.
Whether every such law is the isotropic member of a family asks for an extension off the frame's
shapes, where $\le_\alpha$ relates matrices of different correlation orbits. That question and the
families off the frame's shapes are part of problem~\ref{prob:continuous-fields} of
\S\ref{sec:open-problems}.

\begin{proof}[Proof of Proposition~\ref{prop:diagonal-stable-order-families} (sketch)]
\emph{(1).} Proposition~\ref{prop:similarity-stable-order-families}(1) makes $(\mathrm D_+(d), \le_\alpha)$ a
wedge datum, and the orbits are those of Proposition~\ref{prop:diagonal-ray-families}(1). The relative step
of $D_s^2$ and $D_{s'}^2$ is $D_u^2$, and condition (c) of Proposition~\ref{prop:stable-step-cone}(1) is the
stated inequality squared.

\emph{(2).} (K3) of Definition~\ref{def:gw-cascade-family} with $A = D_s$ gives the frame images. For
$B = \diag(b) \in \mathcal D_\alpha$ and $u_i := e^{tb_i}$, (1) gives $I \le_\alpha D_u^2 = e^{2tB}$, and (K1)
at $0 \le_\alpha I \le_\alpha D_u^2$ is $B$-self-decomposability in the sense of
Proposition~\ref{prop:stable-exponent-cone}; $B = I$ gives self-decomposability.

\emph{(3).} (a) restricts the families of Proposition~\ref{prop:diagonal-ray-families}(3) by
Lemma~\ref{lem:gw-kernels-from-origin}(1), $\le_\alpha$ being contained in the Loewner order. (b) restricts
those of Proposition~\ref{prop:similarity-stable-order-families}(2), the profiles of
Proposition~\ref{prop:orbit-step-cone} included, to $\mathrm D_+(d)$.

\emph{(4).} (a) An increment for $e^{-tB}$ would make $\Theta_r - e^{-tB}\Theta_re^{-tB}$ positive
semidefinite, hence $B\Theta_r + \Theta_rB$, whose principal $2 \times 2$ block has the negative determinant
$4(1+\alpha) - (2+\alpha)^2r^2$. (b) $e^{-|\omega|^\alpha} = \prod_ie^{-|\omega_i|^\alpha}$ fails at
$e_1 + e_2$, the transform being that of Lemma~\ref{lem:elliptical-stable}(2). (c) The convolution of the
two families of (3) is a family by Lemma~\ref{lem:gw-kernels-from-origin}(3); its $\mu_I$ is not rotation
invariant, which the isotropic member of a general linear family is (Lemma~\ref{lem:parameter-space}(5)),
and not a product.
\end{proof}
